\section{Discussion}

Our experiments reveal a critical methodological constraint for gradient subspace analysis rather than evidence for or against the Progressive Gradient Orthogonalization hypothesis.

\subsection{Key Findings}

\textbf{Finding 1: Gradient subspace methods require sufficient accumulation density.}

Accumulating one gradient per epoch produces a subspace whose rank equals the number of epochs, not the requested SVD rank. In our experiment, 10 epochs yielded a rank-10 subspace in a 25-million-parameter space---capturing approximately $10/25\text{M} = 4 \times 10^{-7}$ of variance along any direction.

\textbf{Finding 2: The PGO hypothesis remains unverified, not refuted.}

Our null result stems from measurement failure rather than absence of the hypothesized effect. The question of whether gradient subspaces can capture simplicity bias remains open pending valid measurement.

\textbf{Finding 3: Multi-batch accumulation is necessary.}

With $\sim$38 batches per epoch and 10 accumulation epochs, multi-batch accumulation would yield $\sim$380 gradient samples---sufficient for meaningful rank-50 subspaces.

\subsection{Limitations}

\textbf{The core hypothesis is unverified.} We cannot claim that PGO works or fails because our measurement apparatus prevented valid testing.

\textbf{Only one dataset tested.} Other spurious correlation benchmarks may exhibit different gradient dynamics.

\textbf{No baseline comparison performed.} Comparison against JTT, DFR, or Group DRO awaits valid PGO implementation.

\textbf{Single random seed.} Our experiments used seed 42 for all runs. While sufficient to demonstrate measurement apparatus failure, reproducibility across seeds should be verified in future work.

\subsection{Broader Impact}

This work identifies a methodological pitfall in gradient-based machine learning interventions. Researchers proposing gradient subspace methods should validate accumulation density before claiming directional findings.
